\documentclass[11pt]{amsart}
\usepackage[localtheorems]{raeez-math-template}

\newtheorem{theorem}{Theorem}[section]
\newtheorem{lemma}[theorem]{Lemma}
\newtheorem{proposition}[theorem]{Proposition}
\newtheorem{corollary}[theorem]{Corollary}
\newtheorem{conjecture}[theorem]{Conjecture}
\theoremstyle{definition}
\newtheorem{definition}[theorem]{Definition}
\theoremstyle{remark}
\newtheorem{remark}[theorem]{Remark}

\providecommand{\Ainf}{A_\infty}
\providecommand{\Bmod}{\mathbf{B}_{\mathrm{mod}}}

\title{Curved Dunn additivity at higher genus:\\
the bridge between modular-bootstrap and twisting-cochain complexes}
\author{Raeez Lorgat}
\date{April 2026}

\begin{document}
\maketitle

\begin{abstract}
Classical Dunn additivity decomposes $E_2 \simeq E_1 \otimes E_1$. At
genus $g \ge 1$, the modular bar complex of the curved modular
$\Ainf$-chiral operad is curved by
$d^2_{\mathrm{fib}} = \kappa_{\mathrm{mod}}(A) \cdot \omega_g$ (where
$\omega_g$ is the Arakelov form on $\overline{\cM}_{g,n}$ and
$\kappa_{\mathrm{mod}}(A)$ is the modular curvature coefficient of
the chiral algebra); the Dunn decomposition is obstructed
by $R$-matrix monodromy at sewing circles.

The core obstruction lives in $H^2$ of two different complexes:
a modular-bootstrap complex and a curved-Dunn twisting-cochain
complex. A degree-two comparison map $\Phi$, when supplied, compares
their obstruction classes. Genus~$1$ is a separate curved comparison
governed by the Hodge bundle and the Faltings--Arakelov normalisation
$\omega_1=12c_1(\lambda_1)$.  The higher-genus vanishing statement
starts at $g\ge2$ and is conditional on modular-bootstrap
$H^2$-acyclicity together with the degree-two comparison.  KZB
boundary composition is conditional on compatible boundary formal
types and mixed Stokes data; a simple node is regular singular of
Newton slope~$0$.
\end{abstract}

\section{Genus-stratified obstruction}
\label{sec:three-levels-cd}

\paragraph{Genus $0$ (uncurved Dunn).}
$\cO^{\Ainf\text{-ch}}\big|_{g=0} = \mathrm{SC}^{\mathrm{ch,top}}
\times E_1^{\mathrm{tr}}$. The bar complex factors as an ordinary
tensor product.

\paragraph{Obstruction theory.}
The obstruction to extending the genus-$0$ product to genus $g$
lies in $H^2(\mathrm{TCo}^\bullet_g(A), d_0)$, the twisting-cochain
complex of \S\ref{sec:two-complexes-bridge}. At genus~$1$: controlled
by $\kappa_{\mathrm{mod}}(A) \cdot \omega_1$ via the monodromy
$\mathrm{Mon}(R)$ of the classical $r$-matrix around the $A$-cycle.
At genus $g \ge 2$: conditional on modular-bootstrap $H^2$-acyclicity
and the degree-two comparison map $\Phi$.

\paragraph{Genus-one twisted tensor product.}
\begin{proposition}[Genus-$1$ curved tensor-product comparison]%
\label{prop:genus1-twisted-tensor-product-standalone}%
Let $A$ be a chirally Koszul algebra on a smooth projective curve
$X$ with nonzero modular curvature coefficient
$\kappa_{\mathrm{mod}}(A)$. Assume that the $A$-cycle monodromy of the
classical $r$-matrix defines a twisting cochain $\tau_1$ in the
genus-one Eilenberg--Zilber complex and that
$d\tau_1=\kappa_{\mathrm{mod}}(A)\omega_1$ as a curvature
representative in the chosen coderived ambient category. Then
\[
  \Bmod\bigl(\cO^{\Ainf\text{-ch}}\bigr)\big|_{g=1}
  \;\simeq\;
  \Bmod\bigl(\mathrm{SC}^{\mathrm{ch,top}}\bigr)\big|_{g=1}
  \;\otimes^{\mathrm{Mon}(R)}\;
  \mathbf{B}\bigl(E_1^{\mathrm{tr}}\bigr),
\]
where $\otimes^{\mathrm{Mon}(R)}$ is the Eilenberg--Zilber twisted
tensor product with twisting cochain
$\tau_1 = \mathrm{Mon}(R) = \exp\!\bigl(\oint_{A\text{-cyc}} r(z)\,
dz\bigr)$.
\end{proposition}

\section{Two complexes, one obstruction}
\label{sec:two-complexes-bridge}

Fix a chirally Koszul $A$ on $X$ at non-critical level, modular
characteristic $\kappa_{\mathrm{mod}}(A) \neq 0$, and Arakelov form
$\omega_g$. At genus~$1$ we use the active normalisation
$\omega_1=12c_1(\lambda_1)$.

\begin{definition}[Modular-bootstrap complex]%
\label{def:modular-bootstrap-complex-standalone}%
$\bigl(\mathfrak{g}^A_{\mathrm{mod}},\; d_{\Theta^{(0)}}\bigr)$, the dg
Lie algebra whose underlying graded vector space is
$\prod_{g \ge 0,\, n \ge 0} \mathrm{Hom}_{\Sigma_n}(M\mathrm{Ass}(g,
n),\, \mathrm{End}_A(n))^{\Sigma_n}$, with differential
$d_{\Theta^{(0)}} = [\Theta^{(0)}, -]$ the bracket with the genus-$0$
Maurer--Cartan seed. The genus tower $\{\Theta^{(g)}\}$ satisfies
$D\Theta + \tfrac{1}{2}[\Theta, \Theta] = 0$ genus-wise.
The degree-two vanishing used below is the modular-bootstrap input;
the comparison with the curved-Dunn complex is a separate low-degree
bridge condition.
\end{definition}

\begin{definition}[Curved-Dunn twisting-cochain complex]%
\label{def:curved-dunn-complex-standalone}%
$\bigl(\mathrm{TCo}^\bullet_g(A),\; d_0\bigr)$ at genus $g$ with
\[
  \mathrm{TCo}^\bullet_g(A)
  \;=\; \mathrm{Hom}\!\bigl(
     \mathbf{B}(E_1^{\mathrm{tr}}),\;
     \mathrm{gr}^g\,\Bmod(\mathrm{SC}^{\mathrm{ch,top}})\bigr),
\]
differential $d_0(f) = \partial f - (-1)^{|f|} f \partial$ the
convolution differential of the Eilenberg--Zilber twisted tensor
product. The obstruction to extending
$\tau_{<g}$ to $\tau_g$ is
$\mathrm{Ob}_g \in H^2(\mathrm{TCo}^\bullet_g(A), d_0)$.
\end{definition}

\begin{remark}[Two-complex distinction]%
\label{rem:fm67-two-complex}%
Both complexes host the obstruction $\mathrm{Ob}_g$ but via different
differentials on \emph{a priori} different graded objects:
bootstrap uses the tree-twisted Feynman transform of
$M\mathrm{Ass}$, Dunn uses the Eilenberg--Zilber machine on the
tensor factorisation. The modular-bootstrap $H^2$ vanishing is an
independent input.  The curved-Dunn $H^2$ has no direct argument from
the Arakelov class, since its differential is governed by the
transverse $E_1$ factor. The bridge is the required comparison.
\end{remark}

\section{The bridge proposition}
\label{sec:bridge-cd-standalone}

\begin{proposition}[Modular-bootstrap to curved-Dunn comparison criterion]%
\label{prop:modular-bootstrap-to-curved-dunn-bridge-standalone}%
Assume there is a complete filtered \(L_\infty\)-morphism
\[
  \Phi \;:\;
  \bigl(\mathfrak{g}^A_{\mathrm{mod}},\; d_{\Theta^{(0)}}\bigr)
  \;\longrightarrow\;
  \bigl(\mathrm{TCo}^\bullet_g(A),\; d_0\bigr),
\]
natural in $A$, such that:
\begin{enumerate}[label=\textup{(\roman*)}]
\item the linear term \(\Phi_1\) is strict on the associated graded
  and induces an isomorphism on completed degree-two cohomology;
\item for every $g \ge 1$, the full \(L_\infty\) transport sends the
  bootstrap obstruction class to the curved-Dunn obstruction class
  \(\mathrm{Ob}_g \in H^2(\mathrm{TCo}^\bullet_g(A))\);
\item the inverse-limit comparison has vanishing \(\varprojlim^1\) in
  degree~$2$, so completed cohomology is computed by the associated
  graded comparison.
\end{enumerate}
Under this hypothesis,
$H^2(\mathfrak{g}^A_{\mathrm{mod}}, d_{\Theta^{(0)}}) = 0$ implies
$H^2(\mathrm{TCo}^\bullet_g(A), d_0) = 0$.
\end{proposition}

\begin{proof}
Work in the completed filtered complexes.  Hypotheses \textup{(i)}
and \textup{(iii)} identify the completed degree-two cohomology of
the bootstrap and curved-Dunn complexes through the associated graded
linear term \(\Phi_1\).  The higher Taylor components of \(\Phi\)
transport Maurer--Cartan obstruction classes, so hypothesis
\textup{(ii)} identifies the obstruction class, not a preferred
representative. If
\[
  H^2(\mathfrak{g}^A_{\mathrm{mod}}, d_{\Theta^{(0)}})=0,
\]
then every transported degree-two curved-Dunn class is exact.  In
particular the obstruction class \(\mathrm{Ob}_g\) vanishes in
completed cohomology.
\end{proof}

\begin{remark}[Concrete bridge data]%
\label{rem:bridge-provenance-standalone}%
Both complexes are governed by the same underlying datum:
$R$-matrix monodromy at sewing circles. Modular-bootstrap organises
this datum by stable ribbon graph (Feynman transform); curved-Dunn
organises it by twisting cochain (Eilenberg--Zilber). The candidate
for $\Phi$ sends a stable ribbon graph $\Gamma$ of type $(g,n)$ with
internal edges $e_1,\dots,e_k$ to
\[
  \bigl(\mathrm{Mon}(R)_{e_1}\otimes\cdots\otimes
  \mathrm{Mon}(R)_{e_k}\bigr)\cdot\mathrm{ev}_\Gamma .
\]
The bridge hypothesis is exactly the assertion that this candidate is
a filtered \(L_\infty\) comparison through degree two, identifies the
obstruction classes, and induces an isomorphism on completed $H^2$.
The graph map must carry the automorphism weights and orientation
cocycle; in particular the theta graph has weight \(1/12\) and the
bouquet graph has weight \(1/8\). A
model-categorical route through Bellier-Mill\`es--Drummond-Cole would
have to prove the same degree-two comparison.
\end{remark}

\section{Genus-one curved comparison: Gauss--Manin and Arakelov}
\label{sec:genus1-gauss-manin-standalone}

\begin{proof}[Proof of Proposition~\ref{prop:genus1-twisted-tensor-product-standalone}
under the stated twisting-cochain hypothesis]
The hypothesis supplies the only primitive used below: the
genus-one twisting cochain $\tau_1$ with
$d\tau_1=\kappa_{\mathrm{mod}}(A)\omega_1$ in the selected coderived
model.  The proof verifies that this input gives the displayed
twisted tensor product.

\emph{Step 1 (Gauss--Manin connection on $\Lambda^2 H^1$).}
The universal elliptic curve
$\mathcal{E} \to \overline{\cM}_{1,1}$ carries the Hodge bundle
$\lambda_1 = \pi_* \Omega^1_{\mathcal{E}/\overline{\cM}_{1,1}}$ and
its second wedge
$\lambda_2 = \Lambda^2 R^1\pi_* \mathbb{Q}$. The Gauss--Manin
connection $\nabla^{\mathrm{GM}}$ on
$R^1\pi_* \mathbb{Q}$ descends to $\lambda_2$ as a flat connection
on a rank-$1$ bundle. On the curve side, the bar-complex differential
at genus~$1$ factorises through $\nabla^{\mathrm{GM}}$ via the
fibrewise Hochschild complex:
\[
  d_{\Bmod,\, g=1}
  \;=\;
    d_{\mathrm{fib}} \;+\; \kappa_{\mathrm{mod}}(A) \cdot \omega_1
    \wedge (\cdot),
\]
with $\omega_1=12c_1(\lambda_1)$ in the Faltings--Arakelov
normalisation. The Arakelov class is a base class.  The comparison is
representative-level inside the chosen Eilenberg--Zilber complex.  It
does not assert that the Hodge or Arakelov class is exact on
$\overline{\cM}_{1,1}$.

\emph{Step 2 (Arakelov pairing on $\lambda_1$).}
The Faltings--Arakelov metric on the Hodge bundle fixes the
representative of $\omega_1=12c_1(\lambda_1)$.  We use this
normalisation only as a curvature representative on the base
moduli stack.  The Arakelov pairing
$\langle\cdot, \cdot\rangle_{\mathrm{Ar}}$ on
$\lambda_1$ is $\nabla^{\mathrm{GM}}$-compatible, giving an
isomorphism of flat bundles
$\lambda_2 \otimes \lambda_1 \simeq \lambda_1^\vee$ up to an integer
twist. Pulling back along the chosen genus-one comparison, the
twisted K\"unneth factorisation
becomes
\[
  \Bmod^{(1)}(A)
  \;\simeq\;
  \Bmod^{(1)}(\mathrm{SC}^{\mathrm{ch,top}})
  \otimes \mathbf{B}(E_1^{\mathrm{tr}})
  \quad\text{twisted by}\quad
  \mathrm{Mon}(R) \;=\; \exp\!\bigl(\oint_{A\text{-cyc}} r(z)\, dz\bigr).
\]
The exponent is the $A$-cycle period of the classical $r$-matrix, a
well-defined element of
$\mathrm{End}_A \otimes \mathrm{End}_A$ because $r(z)$ satisfies the
KZ integrability condition.

\emph{Step 3 (Eilenberg--Zilber compatibility).}
The Eilenberg--Zilber twisted tensor product with twisting cochain
$\tau = \mathrm{Mon}(R)$ has differential
$\partial_\tau(x \otimes y) = \partial x \otimes y +
(-1)^{|x|} x \otimes \partial y + \tau \cdot (x \otimes y)$.
The assumption $d\tau_1=\kappa_{\mathrm{mod}}(A)\omega_1$ identifies
this differential with the curved differential at genus~$1$.  The
factorisation is an isomorphism in the stated coderived ambient
category.
\end{proof}

\begin{remark}[Genus-one bridge]%
\label{rem:genus-one-bridge-standalone}%
The genus-one bridge is the first nonzero test case for the
curved-Dunn comparison: Gauss--Manin uncurving on $\Lambda^2H^1$, the
Faltings--Arakelov class, and the Eilenberg--Zilber twisting cochain
must give the same degree-two obstruction class.
\end{remark}

\begin{remark}[Relationship to the bridge at genus~$1$]%
\label{rem:bridge-at-genus-1-standalone}%
At $g = 1$, the bridge comparison $\Phi$ of
Proposition~\ref{prop:modular-bootstrap-to-curved-dunn-bridge-standalone}
reduces to the assertion that the annular bootstrap insertion
$\Theta^{(1)}$ and the curved-Dunn twisting cochain $\tau_1$ represent
the same monodromy class. At $g \ge 2$, the comparison is not
tautological: higher-genus stable-graph strata must match higher
Eilenberg--Zilber corrections.
\end{remark}

\section{Conditional higher-genus curved-Dunn $H^2$ vanishing}
\label{sec:higher-genus-vanishing-standalone}

\begin{theorem}[Conditional curved-Dunn $H^2$ vanishing for $g\ge2$]%
\label{thm:curved-dunn-H2-vanishing-higher-genus-standalone}%
Assume the modular-bootstrap degree-two vanishing and the low-degree
bridge hypothesis of
Proposition~\ref{prop:modular-bootstrap-to-curved-dunn-bridge-standalone}.
For every $g \ge 2$ and every chirally Koszul $A$ in the standard
landscape at non-critical level,
\[
  H^2\bigl(\mathrm{TCo}^\bullet_g(A),\; d_0\bigr) \;=\; 0.
\]
Consequently, $\tau_{<g}$ extends to $\tau_g$ for every $g\ge2$ up to
convolution-gauge equivalence, and the modular bar complex factorises
cohomologically as a twisted K\"unneth tensor product:
\[
  \Bmod\bigl(\cO^{\Ainf\text{-ch}}\bigr)
  \;\simeq\;
  \Bmod\bigl(\mathrm{SC}^{\mathrm{ch,top}}\bigr)
  \;\otimes^{\tau}\;
  \mathbf{B}\bigl(E_1^{\mathrm{tr}}\bigr).
\]
\end{theorem}

\begin{proof}
$H^2(\mathfrak{g}^A_{\mathrm{mod}}, d_{\Theta^{(0)}}) = 0$ at every
$g \ge 2$ by the modular-bootstrap input.  The genus-one curved
comparison of Proposition~\ref{prop:genus1-twisted-tensor-product-standalone}
does not enter as a vanishing base case.  The Maurer--Cartan recursion
exhibits the obstruction class; it does not by itself prove
vanishing.

Apply Proposition~\ref{prop:modular-bootstrap-to-curved-dunn-bridge-standalone}(iii):
the assumed isomorphism
$H^2(\mathfrak{g}^A_{\mathrm{mod}}) \simeq
H^2(\mathrm{TCo}^\bullet_g(A))$
transports the vanishing. The extension
$\tau_{<g} \mapsto \tau_g$ is automatic:
$\mathrm{Ob}_g = [d_0(\tau_{<g}) +
\tfrac{1}{2}[\tau_{<g}, \tau_{<g}]_{\mathrm{conv}}] = 0$ in
$H^2$, so the cocycle is a coboundary and the twisting cochain
extends.
\end{proof}

\section{Boundary KZB formal type and Stokes composition}
\label{sec:jimbo-miwa-standalone}

At non-critical level, smooth-locus KZB flatness does not by itself
define irregular modular-operad composition at the boundary. The needed
input is a compatible meromorphic formal type on every stable-graph
boundary stratum, together with mixed-channel Stokes data at corners.
Edgewise Stokes data alone do not determine two-edge clutching.

Fix a point $p \in \partial\overline{\cM}_{g, n}$ on the boundary
nodal divisor, with local coordinate $\hbar_p$ transverse to the
divisor. The KZB connection near $p$ has the form
$\nabla^{\mathrm{KZB}} = d + \Omega^{\mathrm{KZB}}(\hbar_p)$. At a
simple node the polar part is logarithmic:
$\Omega^{\mathrm{KZB}}(\hbar_p)=A_1\,d\hbar_p/\hbar_p+
\Omega_{\mathrm{reg}}$. The singularity is regular, with Newton
slope~$0$; it is not a Stokes source. Genuine irregular Stokes data
enter only on strata whose formal boundary model contains higher-order
terms after ramification.

\begin{theorem}[Stokes-sector composition from boundary KZB formal type]%
\label{thm:irregular-singular-kzb-regularity-standalone}%
Assume the boundary formal-type theorem: on every stable-graph
boundary stratum, including two-edge corners, the KZB connection has
compatible formal types, regular-singular slope-$0$ simple-node
restrictions, residue-gap nonresonance in the irregular sectors, and
vanishing mixed-channel Stokes curvature. Then for every $(g,n)$ at
non-critical level, the resulting formal/Stokes boundary data define a
factorisation map on the modular operad $M\mathrm{Ass}(g,n)$.
Consequently, modular-operad composition in the irregular
boundary-formal-type regime is defined for all $(g,n)$.
\end{theorem}

\begin{proof}[Proof under the boundary formal-type hypothesis]
At a simple node the logarithmic residue supplies regular-singular
monodromy and no Stokes filtration. On every boundary component with
nonzero irregular slope, the assumed formal type gives the exponential
factors, Stokes directions, and Stokes groups after finite
ramification. Jimbo--Miwa isomonodromic deformation and the
formal/Stokes classification of meromorphic connections identify the
sectorial continuation maps of flat sections; the Stokes matrices are
additional monodromy data constrained by the formal type and by the
isomonodromy equations.

These one-edge operations do not imply coherence at a two-edge corner.
The assumed mixed-channel zero-curvature condition is exactly the
statement that the two orders of clutching yield the same sectorial
continuation. With this corner compatibility imposed on every
stable-graph stratum, Getzler--Kapranov clutching composes the
regular-singular simple-node monodromies and the nonzero-slope Stokes
sector maps into a well-defined modular-operad composition.
\end{proof}

\begin{corollary}[Modular operad composition under boundary formal type]%
\label{cor:modular-operad-composition-all-levels-standalone}%
Under the boundary formal-type theorem, the modular
$\Ainf$-chiral operad $\cO^{\Ainf\text{-ch}}$ admits
operadic composition at integer levels (classical regular-singular
KZB) and at non-critical levels in the supplied formal-type sectors
(regular-singular simple-node monodromy plus Stokes-sector
composition on nonzero-slope boundary strata), uniformly at all
$(g,n)$ satisfying the hypothesis.
\end{corollary}

\section{Summary}
\label{sec:cd-summary-standalone}

\begin{enumerate}[label=\textup{(\arabic*)}]
\item \textbf{Bridge.}
  Proposition~\ref{prop:modular-bootstrap-to-curved-dunn-bridge-standalone}
  states the degree-two comparison criterion between the
  modular-bootstrap and curved-Dunn obstruction complexes.
\item \textbf{Genus-$1$ curved comparison.}
  Proposition~\ref{prop:genus1-twisted-tensor-product-standalone} uses
  Gauss--Manin on $\Lambda^2 H^1$ and Arakelov pairing on
  $\lambda_1$ under the stated twisting-cochain hypothesis.
\item \textbf{Higher-genus vanishing.}
  Theorem~\ref{thm:curved-dunn-H2-vanishing-higher-genus-standalone}: curved-Dunn
  $H^2 = 0$ for $g \ge 2$ under modular-bootstrap acyclicity and the
  degree-two bridge hypothesis; the twisting cochain extends up to
  gauge and the twisted tensor factorisation holds cohomologically.
\item \textbf{Boundary KZB formal type.}
  Theorem~\ref{thm:irregular-singular-kzb-regularity-standalone}: modular-operad
  composition in the irregular boundary-formal-type sector under
  compatible formal types, residue-gap nonresonance, and mixed Stokes
  data.
\end{enumerate}

Under these hypotheses, the curved-Dunn decomposition has the
cohomological, gauge-qualified form stated above: the modular bar
complex is modeled by a twisted Eilenberg--Zilber tensor product of
the closed-colour $\mathrm{SC}^{\mathrm{ch,top}}$ modular bar with
the open-colour transverse $E_1$ bar, with twisting cochain built
recursively from $R$-matrix monodromy at genus-raising clutching.

\begin{thebibliography}{99}
\bibitem{Vol1} R.~Lorgat,
\emph{Chiral bar and chiral cobar: Koszul duality, modular invariants,
and the quantum geometry of vertex algebras}, Volume~I (2026).
\bibitem{Vol2} R.~Lorgat,
\emph{$A_\infty$-chiral algebras and 3D holomorphic-topological QFT},
Volume~II (2026).
\bibitem{JimboMiwa81} M.~Jimbo and T.~Miwa,
\emph{Monodromy preserving deformation of linear ordinary
differential equations with rational coefficients. II},
Physica D~2 (1981).
\bibitem{BMDC} J.~Bellier-Mill\`es and G.~Drummond-Cole,
\emph{Model structures on curved operads} (preprint, 2024).
\end{thebibliography}

\end{document}
